\documentclass[12pt]{article}
\usepackage[utf8]{inputenc}
\usepackage[T1]{fontenc}
\usepackage{hyperref}
\title{Mangerr Language Resources for K-12 Curriculum Development}
\author{Compiled by Claude Sonnet 5.5 (claude-sonnet-5-5) for Andr\'as Kornai}
\date{}
\begin{document}
\maketitle

\section{Introduction}
Mangerr (Mengerrdji; ISO 639-3 zme; Glottocode mang1382) is one of three varieties of the extinct Giimbiyu language of the Alligator Rivers region of the Northern Territory, Australia. The other two are Urningangga (Wuningak) and Erri (Arri). Sources give the end of the language as sometime between the 1980s and 1990s, and ASJP lists it as extinct since 1990. The name Giimbiyu is a Gaagudju word meaning ``of the stoney country''. Descendants of the Giimbiyu people adopted Kunwinjku. The searches found no school programme for the language.

\section{Primary Educational Resources}
No primers, textbooks or school programmes were found in the searches.

\section{Linguistic and Reference Materials}
\begin{itemize}
\item Bruce Birch (2006), \emph{A first dictionary of Erre, Mengerrdji and Urningangk}. The title appeared in a Wikipedia summary. I did not find a direct link to the dictionary itself.
\item Wikipedia, Giimbiyu language: \url{https://en.wikipedia.org/wiki/Giimbiyu_language}
\item Wikipedia, ISO 639 redirect for zme: \url{https://wikipedia.com/wiki/ISO_639:zme}
\item Glottolog entry for the Giimbiyu family: \url{https://glottolog.org/resource/languoid/id/urni1238}
\item ASJP wordlist entry for Mangerr: \url{https://asjp.clld.org/languages/MANGERR}
\end{itemize}

\section{Religious and Community Resources}
\begin{itemize}
\item Wikipedia, Giimbiyu people (background on the people and their descendants): \url{https://en.wikipedia.org/wiki/Giimbiyu_people}
\end{itemize}

\section{Digital and Audio Resources}
No audio or digital teaching resources were found in the searches.

\section{Curriculum Development Context}
The language has no living speakers, and the documentation visible online is limited to a small wordlist and a dictionary of the three varieties. A revitalization curriculum would build on Birch's dictionary and archival material, and would need the agreement of the descendant communities, who now mainly speak Kunwinjku.

\end{document}